% !TeX program = pdflatex
% Draft of the block "The Carried Insertions at Complex Data" for WILSON_LOOPS_CONTINUUM_LIMIT.tex, Section 6.
% Written by a second session, which edits only this file and COMPLEX_COVARIANCE_STATE.md; the session of the
% rewrite of the inductive step moves the section into the paper. State, sources, conditions, and what the
% integrating session has to match: COMPLEX_COVARIANCE_STATE.md.
% Everything between \begin{document} and \end{document} below the title is the text for the paper; the
% preamble is the paper's, with xr for the paper's labels (compile the paper first in the same directory).
\documentclass[11pt]{preprint}
\usepackage[T1]{fontenc}
\usepackage[utf8]{inputenc}
\usepackage{amsmath,amssymb,amsthm}
\usepackage{mathtools}
\usepackage{mathrsfs}
\usepackage[a4paper,margin=28mm]{geometry}
\usepackage{enumitem}
\usepackage{graphicx}
\usepackage{xr-hyper}
\usepackage[hidelinks]{hyperref}
\externaldocument[][nocite]{WILSON_LOOPS_CONTINUUM_LIMIT}[../Continuum Limit Manuscript Draft/WILSON_LOOPS_CONTINUUM_LIMIT.pdf]
\newcommand{\dd}{\mathrm{d}}
\newcommand{\id}{\mathbf 1}
\newcommand{\diam}{\operatorname{diam}}
\newcommand{\dist}{\operatorname{dist}}
\DeclarePairedDelimiter{\abs}{\lvert}{\rvert}
\DeclarePairedDelimiter{\norm}{\lVert}{\rVert}
\setlength{\parskip}{0.6em}
\emergencystretch=3em
\title{Draft: The Carried Insertions at Complex Data}
\author{}
\date{}
\numberwithin{equation}{section}
\newtheorem{theorem}{Theorem}[section]
\newtheorem{proposition}[theorem]{Proposition}
\newtheorem{lemma}[theorem]{Lemma}
\newtheorem{corollary}[theorem]{Corollary}
\theoremstyle{definition}
\newtheorem{definition}[theorem]{Definition}
\theoremstyle{remark}
\newtheorem{remark}[theorem]{Remark}
\theoremstyle{plain}
\renewcommand{\qedsymbol}{\textit{q.e.d.}}
\begin{document}
\maketitle

\section{The Carried Insertions at Complex Data}\label{sec:carried-complex}

The clause \eqref{eq:sup-clause} of Definition~\ref{def:dressed-loop} bounds the insertions on the whole space of the step, and not only at its real backgrounds. For the groups which a step makes this is part of what the count of that step delivers, in the form \eqref{eq:production-clause}. For a group $W$ which was made at an earlier step $j$ and is carried to the step $k + 1$ the question is of another kind, since its insertion is the function which was bounded at the step $j$, and it has now to be bounded on the space of the step $k + 1$ and summed over the groups at a traversal of that step. The obvious bound fails. In the units of the fine lattice the space of the step $k + 1$ lies in that of the step $j$, so that the supremum $\norm{Y_W}_{(j)}$ of the step $j$ bounds the insertion; but a traversal of the step $k + 1$ consists of $L^m$ traversals of the step $j$, $m \coloneqq k + 1 - j$, each with carried groups of its own, and the sum of these suprema is larger than the bound of the step $j$ by the factor $L^m$. At the real backgrounds this factor is compensated by the contraction of the deviation (Section~\ref{sec:contraction}). At complex data there is no deviation to contract, and the compensation has to come from the size of the complex part: the radii of the space of the step $k + 1$ are smaller than those of the step $j$ by the factor $L^{-m}$, up to the ratio of the couplings, so that Schwarz's lemma on the larger space gains $L^{-m}$ for each insertion. This cancels the number of the traversals and no more, and the sum over the generations $m = 1, \dots, k$ would still grow with $k$.

What removes the growth is gauge covariance at complex gauge transformations. The complex part of a configuration of the space of the step $k + 1$ varies slowly, by the gradient condition of that space, and on a set which is small in the units of the scale $k + 1$ it is therefore, to first order, a complex gauge transformation with a gauge function linear in the position. The dependence set of a carried group is such a set, of diameter $\delta_W$ of the order of $ML^{-m}\sigma(W)$ in these units. Removing the linear part by a gauge transformation which is the identity at the first site of the interval of $W$ leaves a complex part smaller by the factor $\delta_W$, and Schwarz's lemma then gains $L^{-2m}\sigma(W)$, which is summable over the generations; the factor $\sigma(W)$, and the groups whose dependence sets are not small, are paid by a sixteenth of the weight $e^{\nu\sigma(W)}$. This is the reason for the rate $\lambda_c < 1$ in \eqref{eq:sup-clause}. The insertions are covariant under gauge transformations with values in the compact group only, and the first step is to see that the covariance extends.

Throughout this section we work on the labelled term without large-field regions, we fix steps $1 \leq j \leq k$ and put $m \coloneqq k + 1 - j$, and we measure lengths in the units of the fine lattice, so that a traversal of the step $n$ consists of $L^n$ fine bonds and a cube of $\pi_n$ has the side $ML^n$. We write $\mathfrak g$ for the Lie algebra of $G$, $\mathfrak g^c \coloneqq \mathfrak g \otimes \mathbb C$, and $G^c$ for the complexification of $G$, the connected complex Lie group $G\exp(i\mathfrak g)$ of matrices with the Lie algebra $\mathfrak g^c$; a representation $\pi$ of $G$ extends to a holomorphic representation of $G^c$. The norm $\abs{\,\cdot\,}$ is that of Definition~\ref{def:dressed-loop}, on the matrices of the size of $G$, and it is invariant under the conjugation by the elements of $G$.

\subsection{The Spaces of the Steps in Fine Units}\label{sec:spaces-fine-units}

For a configuration $U$ with values in $G$ and a field $\alpha$ on the fine bonds with values in $\mathfrak g^c$ we write $e^{i\alpha}U$ for the configuration $b \mapsto e^{i\alpha(b)}U(b)$, with values in $G^c$. For a bond $b = \langle x, x + e_\mu\rangle$ and a unit vector $e$, one of the $\pm e_\nu$, let $b + e \coloneqq \langle x + e, x + e + e_\mu\rangle$, and let the \emph{covariant difference} of $\alpha$ be
\begin{equation}\label{eq:covariant-difference}
    (\delta_e\alpha)(b) \coloneqq U(x, x + e)\,\alpha(b + e)\,U(x, x + e)^{-1} - \alpha(b)\;,
\end{equation}
where $U(x, x + e)$ is the value of $U$ on the bond from $x$ to $x + e$. For $n \geq 2$ let $\mathcal S_n(U)$ be the set of the pairs $(e^{i\alpha}U, \mathbf J)$, with fields $\alpha$ and $\mathbf J$ on the fine bonds with values in $\mathfrak g^c$, such that for all bonds $b$, unit vectors $e$, and fine plaquettes $p$
\begin{equation}\label{eq:space-of-step}
\begin{gathered}
    \abs{\alpha(b)} < \alpha_{1,n}L^{-n}\;,
    \qquad
    \abs{(\delta_e\alpha)(b)} < \alpha_{1,n}L^{-2n}\;,\\
    \abs{(e^{i\alpha}U)(\partial p) - 1} < c_F\,\alpha_{0,n}L^{-2n}\;,
    \qquad
    \abs{\mathbf J(b)} < \alpha_{0,n}L^{-3n}\;,
\end{gathered}
\end{equation}
with a constant $c_F \in (0, 1]$ of the class; $\mathcal S_1(U)$ is the set of the configurations $e^{i\alpha}U$ with $\abs{\alpha(b)} < \alpha_{1,1}L^{-1}$, without a variable $\mathbf J$. These are the conditions of Definition~\ref{def:dressed-loop}(iii) on the term without large-field regions. Indeed, in the units of the scale $n$ the fine lattice has the spacing $\xi = L^{-n}$, the complex part is written $e^{i\xi A'}$, so that $\alpha = \xi A'$, and the conditions \cite[(2.39), (2.34), (2.36), p.~261]{BalabanCRE} on the top domain read $\abs{A'} < \alpha_{1,n}$, $\abs{\nabla^\xi_UA'} < \alpha_{1,n}$ with $\nabla^\xi_U = \xi^{-1}\delta$, $\abs{\partial\mathbf U - 1} < \alpha_{0,n}\xi^2$, and $\abs{\mathbf J} < \alpha_{0,n}$; a current is of the dimension of the third power of an inverse length \cite[(1.8), p.~261]{BalabanRG1}, whence the factor $L^{-3n}$ in fine units. For a set $X$ of bonds we write $\mathcal S_n(U)|_X$ for the set defined by the same conditions on the bonds of $X$ and on the plaquettes and differences which they contain. This set is open in the space of the pairs of a $G^c$-valued configuration and a $\mathfrak g^c$-valued field on $X$, because $\alpha \mapsto e^{i\alpha}U(b)$ maps a neighbourhood of $0$ in $\mathfrak g^c$ biholomorphically onto a neighbourhood of $U(b)$ in $G^c$, and a gauge transformation $v$ with values in $G$ maps $\mathcal S_n(U)|_X$ onto $\mathcal S_n(U^v)|_X$, by the invariance of the norm.

We use the following properties of the class, for a group $W$ of the deviation species made at the step $j$; it has one interval $I_W$, from a site $x_W$ to a site $x'_W$, which lies in its domain $X_W$.
\begin{enumerate}[label=(\alph*)]
\item The insertion $Y_W$ is a function of the restriction of $(\mathbf U, \mathbf J)$ to the dependence set $X^+_W$ (of $\mathbf U$ alone for $j = 1$). It is defined and analytic on the open set
\begin{equation}\label{eq:production-space}
    \mathcal D_W \coloneqq \bigcup_U\mathcal S_j(U)|_{X^+_W}\;,
\end{equation}
the union over the configurations $U$ which are the trivial configuration, the background $U_j(V_j)$ of a block field $V_j$ in the support of the test functions of the term, or a gauge transform of one of these by a $G$-valued gauge transformation. On $\mathcal D_W$ it is bounded by $\norm{Y_W}_{(j)}$, and for $G$-valued gauge transformations $v$
\begin{equation}\label{eq:covariance-law}
    Y_W(\mathbf U^v, R(v)\mathbf J) = \pi(v(x_W))\,Y_W(\mathbf U, \mathbf J)\,\pi(v(x'_W))^{-1}\;,
\end{equation}
where $\mathbf U^v(b) = v(b_-)\mathbf U(b)v(b_+)^{-1}$ and $(R(v)\mathbf J)(b) = v(b_-)\mathbf J(b)v(b_-)^{-1}$ \cite[(1.10), p.~262]{BalabanRG1}, and $\pi$ is the representation of the loop of $W$ (Definition~\ref{def:dressed-loop}(iii) and (vii)). At a later step the insertion of $W$ is the same function of the configuration and of the current, both in fine units.
\item For $V_{k+1}$ in the support of the step $k + 1$ the background $U_{k+1}$ is a background of the step $j$: with $V^*_i \coloneqq M^i(U_{k+1})$ for $j \leq i \leq k$, the block field $V^*_i$ lies in the support of the step $i$ (Supplement, Proposition~S.8, by descending induction on $i$), and $U_i(V^*_i) = U_{k+1}$ by the uniqueness of the minimal orbit (Section~\ref{sec:evaluation-point}). Moreover $\abs{U_{k+1}(\partial p) - 1} < 2B_3\varepsilon_{k+1}L^{-2(k+1)}$ for every fine plaquette $p$ \cite[p.~268]{BalabanCRE}, and the pair of $U_{k+1}$ and its current $J_{k+1}$ is a point of $\mathcal S_{k+1}(U_{k+1})$.
\item The domain $X_W$ is a connected union of at most $u_\star\sigma(W)$ cubes of $\pi_j$, and $I_W$ has at most $u_\star M\sigma(W)L^j$ fine bonds (Definition~\ref{def:dressed-loop}(vii)).
\item The radii of two steps $n' < n$ satisfy $\alpha_{i,n} \leq \varkappa_{n-n'}\,\alpha_{i,n'}$ for $i = 0, 1$, with $\varkappa_m \coloneqq (1 + \beta_0)^2m^{1/2}$, by \eqref{eq:tube-two-six} and \eqref{eq:tube-two-seven}, and $\alpha_{1,n} = \lambda_\alpha\alpha_{0,n}$ with $\lambda_\alpha \coloneqq C^\alpha_1/C^\alpha_0$, since $q_1 = q_0$ (Appendix~\ref{app:constants}).
\end{enumerate}
We assume
\begin{equation}\label{eq:carried-blocking-condition}
    L \geq 20(1 + \beta_0)^2\;,
\end{equation}
so that $\varkappa_mL^{-m} \leq (1 + \beta_0)^2L^{-1} \leq \frac{1}{20}$ for every $m \geq 1$, because $m^{1/2}L^{-(m-1)} \leq 1$ for $L \geq 2$.

\begin{lemma}[The Space of a Later Step Lies in That of an Earlier One]\label{lem:space-inclusion}
Let $U$ be a $G$-valued configuration. Then $\mathcal S_{k+1}(U) \subset \mathcal S_j(U)$.
\end{lemma}

\begin{proof}
By (d) and \eqref{eq:carried-blocking-condition}, $\alpha_{i,k+1}L^{-(k+1)} \leq \varkappa_mL^{-m}\alpha_{i,j}L^{-j} \leq \alpha_{i,j}L^{-j}$, and the same holds with the powers $L^{-2(k+1)}$ and $L^{-3(k+1)}$, which are smaller. Each of the four radii of \eqref{eq:space-of-step} at $n = k + 1$ is therefore at most the corresponding radius at $n = j$, and for $j = 1$ only the first condition is posed.
\end{proof}

\subsection{Complex Gauge Transformations}\label{sec:complex-gauge}

\begin{lemma}[Complexified Covariance]\label{lem:complexified-covariance}
Let $X$ be a finite set of bonds with the set $S$ of its sites, $x, x' \in S$, $\mathcal D$ an open set of pairs $q = (\mathbf U, \mathbf J)$ of a $G^c$-valued configuration and a $\mathfrak g^c$-valued field on $X$ which is invariant under the $G$-valued gauge transformations, and $F$ a holomorphic function on $\mathcal D$ with values in the linear maps of the space of a representation $\pi$ of $G$, such that $F(q^v) = \pi(v(x))F(q)\pi(v(x'))^{-1}$ for every $q \in \mathcal D$ and every $v \in G^S$, where $q^v \coloneqq (\mathbf U^v, R(v)\mathbf J)$. Let $q \in \mathcal D$ and let $t \mapsto u_t$, $t \in [0, 1]$, be a continuous path in $(G^c)^S$ with $u_0 = 1$ and $q^{u_t} \in \mathcal D$ for every $t$. Then
\begin{equation}\label{eq:complexified-covariance}
    F(q^{u_1}) = \pi(u_1(x))\,F(q)\,\pi(u_1(x'))^{-1}\;.
\end{equation}
\end{lemma}

\begin{proof}
The map $u \mapsto q^u$ is holomorphic on $(G^c)^S$, so that $\mathcal U \coloneqq \{ u \in (G^c)^S \,|\, q^u \in \mathcal D \}$ is open, and it contains $G^S$ by the invariance of $\mathcal D$. Since $\pi$ is holomorphic on $G^c$, the function
\[
    \Psi(u) \coloneqq \pi(u(x))^{-1}F(q^u)\pi(u(x')) - F(q)
\]
is holomorphic on $\mathcal U$, and it vanishes on $G^S$ by hypothesis. In the chart $Z \mapsto \exp(Z)$ of $(G^c)^S$ at the identity, with $Z \in (\mathfrak g^c)^S = \mathfrak g^S \oplus i\mathfrak g^S$, the function $\Psi \circ \exp$ is holomorphic near $0$ and vanishes on the real subspace $\mathfrak g^S$, whose complex span is the whole space. All its derivatives at $0$ in the real directions vanish, hence, by complex linearity of the derivatives of a holomorphic function, all its derivatives at $0$, and $\Psi$ vanishes on a neighbourhood of the identity. By the identity theorem it vanishes on the connected component of $\mathcal U$ which contains the identity, and the path $u_t$ lies in this component, because it lies in $\mathcal U$ and starts at the identity. Hence $\Psi(u_1) = 0$, which is \eqref{eq:complexified-covariance}.
\end{proof}

We use the lemma for the insertion of $W$ with the transport along its interval removed. Let $\mathbf U(I_W)$ be the ordered product of $\mathbf U$ along the fine bonds of $I_W$, from $x_W$ to $x'_W$, an element of $G^c$, and
\begin{equation}\label{eq:base-point-form}
    \hat Y_W(\mathbf U, \mathbf J) \coloneqq Y_W(\mathbf U, \mathbf J)\,\pi\bigl(\mathbf U(I_W)\bigr)^{-1}\;.
\end{equation}
Since $\mathbf U^v(I_W) = v(x_W)\mathbf U(I_W)v(x'_W)^{-1}$, the function $\hat Y_W$ is holomorphic on $\mathcal D_W$ and satisfies \eqref{eq:covariance-law} with $x_W$ in the place of $x'_W$. By Lemma~\ref{lem:complexified-covariance}, for a path $u_t$ as there with $u_1(x_W) = 1$,
\begin{equation}\label{eq:base-point-invariance}
    \hat Y_W(\mathbf U^{u_1}, R(u_1)\mathbf J) = \hat Y_W(\mathbf U, \mathbf J)\;.
\end{equation}

The gauge transformation which we shall use is built from the following function. Three properties of the norm enter, each with a constant which depends only on $G$ and on the norm; we fix $\rho_0 \in (0, 1]$ and $c_G \geq 1$ such that all three hold.
\begin{enumerate}[label=(N\arabic*)]
\item For $X, Q \in \mathfrak g^c$ with $\abs X \leq \rho_0$, $\abs{e^{iX}Qe^{-iX} - Q} \leq c_G\abs X\abs Q$.
\item For $X, Y, Z, A \in \mathfrak g^c$ of norm less than $\rho_0$ the matrix $e^{iX}e^{iY}e^{iZ}e^{-i(X + A)}e^{-iZ}$ is within the radius of convergence of the logarithm about $1$, and
\begin{equation}\label{eq:reduction-function}
    \Phi(X, Y, Z, A) \coloneqq -i\log\bigl(e^{iX}e^{iY}e^{iZ}e^{-i(X + A)}e^{-iZ}\bigr)
\end{equation}
is a holomorphic function with values in $\mathfrak g^c$; the function $\Psi \coloneqq \Phi - (Y - A)$ has all its first and second partial derivatives bounded by $c_G$ in norm where the four arguments have norm at most $\rho_0/2$.
\item Let $U$ be a $G$-valued configuration and $\alpha$ a field with $\abs{\alpha(b)} \leq a_\circ \leq \rho_0$ and $\abs{(\delta_e\alpha)(b)} \leq d_\circ \leq a_\circ$ on the bonds of a fine plaquette $p$. Then
\begin{equation}\label{eq:complex-curvature}
    \abs{(e^{i\alpha}U)(\partial p) - 1} \leq c_G\bigl(\abs{U(\partial p) - 1} + d_\circ + a_\circ^2\bigr)\;.
\end{equation}
\end{enumerate}
Such constants exist. (N1) holds because $Q \mapsto e^{iX}Qe^{-iX} - Q$ is holomorphic in $X$ and vanishes at $X = 0$. In (N2) the product is a holomorphic function of the arguments which equals $1$ at the origin, the logarithm of a product of exponentials of small elements of $\mathfrak g^c$ is given by the series of Baker, Campbell, and Hausdorff and lies in $\mathfrak g^c$, and the derivatives are bounded by Cauchy's estimates. For (N3), let $b_1 = \langle x, x + e_\mu\rangle$, $b_2$, $b_3$, $b_4 = \langle x, x + e_\nu\rangle$ be the bonds of $p$ in the order $(\partial p) = b_1b_2b_3^{-1}b_4^{-1}$, $U_i \coloneqq U(b_i)$, $\alpha_i \coloneqq \alpha(b_i)$, and $F \coloneqq U(\partial p)$. Moving the factors $U_i$ to the right we find
\begin{gather*}
    (e^{i\alpha}U)(\partial p) = e^{i\alpha_1}e^{i\beta_2}e^{-i\beta_3}e^{-i\beta_4}F\;,\\
    \beta_2 \coloneqq U_1\alpha_2U_1^{-1}\;,
    \qquad
    \beta_3 \coloneqq FU_4\alpha_3U_4^{-1}F^{-1}\;,
    \qquad
    \beta_4 \coloneqq F\alpha_4F^{-1}\;.
\end{gather*}
Here $U_1\alpha_2U_1^{-1} = \alpha_4 + (\delta_{e_\mu}\alpha)(b_4)$ and $U_4\alpha_3U_4^{-1} = \alpha_1 + (\delta_{e_\nu}\alpha)(b_1)$ by \eqref{eq:covariant-difference}, and the conjugation by $F$ changes an element $Q$ by at most a constant times $\abs{F - 1}\abs Q$. The sum $\alpha_1 + \beta_2 - \beta_3 - \beta_4$ of the exponents is therefore at most $2d_\circ$ plus a constant times $\abs{F - 1}a_\circ$ in norm, the product of the four exponentials differs from $1 + i(\alpha_1 + \beta_2 - \beta_3 - \beta_4)$ by at most a constant times $a_\circ^2$, since its arguments are bounded by $2a_\circ$, and the multiplication by $F$ changes norms by at most a constant factor. This gives \eqref{eq:complex-curvature}.

Two consequences of the definition \eqref{eq:reduction-function} will be used: $\Phi(X, 0, 0, 0) = 0$ for every $X$ and $\Phi(0, 0, Z, 0) = 0$ for every $Z$, and the linear part of $\Phi$ at the origin is $Y - A$. Hence, where the four arguments have norm at most $\rho_0/2$,
\begin{equation}\label{eq:reduction-remainder}
\begin{gathered}
    \abs{\Psi(X, Y, Z, A)} \leq c_G\bigl[(\abs Y + \abs A)(\abs X + \abs Y + \abs Z + \abs A) + \abs X\abs Z\bigr]\;,\\
    \abs{\partial_X\Psi} \leq c_G(\abs Y + \abs Z + \abs A)\;,
    \qquad
    \abs{\partial_Z\Psi} \leq c_G(\abs X + \abs Y + \abs A)\;,\\
    \abs{\partial_Y\Psi} \leq c_G(\abs X + \abs Y + \abs Z + \abs A)\;.
\end{gathered}
\end{equation}
For the bounds of the derivatives we note that $\partial_X\Psi$ vanishes at $(X, 0, 0, 0)$, $\partial_Z\Psi$ at $(0, 0, Z, 0)$, and $\partial_Y\Psi$ at the origin, and we integrate the second derivatives along the segment from these points. For the first line, $\Psi(X, 0, Z, 0)$ vanishes if $X = 0$ or $Z = 0$, so that it equals the integral of its mixed second derivative over $[0, 1]^2$ along $(sX, tZ)$ and is at most $c_G\abs X\abs Z$; and $\Psi(X, Y, Z, A) - \Psi(X, 0, Z, 0)$ is the integral of the derivative in $(Y, A)$ along $(X, sY, Z, sA)$, which vanishes at the origin and is therefore at most $c_G(\abs X + \abs Y + \abs Z + \abs A)$.

\begin{lemma}[The Gauge Reduction of a Slowly Varying Complex Part]\label{lem:gauge-reduction}
Let $B$ be a cube of the fine lattice of side $D \geq 1$ and $x_0$ a site of $B$ with $x_0 + e_\mu \in B$ for every $\mu$. Let $U = e^{i\zeta}$ on the bonds of $B$, with a field $\zeta$ with values in $\mathfrak g$, and let $\alpha$ be a field on the bonds of $B$ with values in $\mathfrak g^c$, such that for all bonds $b$ and unit vectors $e$ with $b, b + e$ in $B$
\begin{equation}\label{eq:reduction-hypotheses}
    \abs{\alpha(b)} \leq a\;,
    \qquad
    \abs{(\delta_e\alpha)(b)} \leq a'\;,
    \qquad
    \abs{\zeta(b)} \leq z\;,
    \qquad
    \abs{\zeta(b + e) - \zeta(b)} \leq z'\;,
\end{equation}
with $a' \leq a$ and $4aD, a, z \leq \rho_0/2$. Put
\begin{equation}\label{eq:gauge-function}
    \omega(x) \coloneqq \sum_\mu a_\mu\,(x - x_0)_\mu\;,
    \qquad
    a_\mu \coloneqq \alpha(\langle x_0, x_0 + e_\mu\rangle)\;,
    \qquad
    u_t \coloneqq e^{it\omega}\;,
\end{equation}
for the sites $x$ of $B$ and $t \in [0, 1]$. Then $u_t(x_0) = 1$, $\abs{\omega} \leq 4aD$, and $(e^{i\alpha}U)^{u_t} = e^{i\alpha_t}U$ on the bonds of $B$, with fields $\alpha_t$ with values in $\mathfrak g^c$ which satisfy, with $a'' \coloneqq a' + c_Gza$ and $r \coloneqq c_G[2a(4aD + 2a + z) + 4aDz]$,
\begin{equation}\label{eq:reduced-phase}
    \abs{\alpha_1(b)} \leq 4Da'' + r\;,
    \qquad
    \abs{\alpha_t(b)} \leq 2a + r\;,
\end{equation}
\begin{equation}\label{eq:reduced-difference}
    \abs{(\delta_e\alpha_t)(b)} \leq a'' + c_Gz(2a + r) + c_G\bigl[(2a + z)a + (4aD + 2a + z)a'' + (4aD + 2a)z'\bigr]\;.
\end{equation}
\end{lemma}

\begin{proof}
The value $\omega(x_0) = 0$ and the bound on $\omega$ are immediate from \eqref{eq:gauge-function}, since $\abs{a_\mu} \leq a$ and $\sum_\mu\abs{(x - x_0)_\mu} \leq 4D$. For a bond $b = \langle x, x + e_\mu\rangle$ of $B$ we have $\omega(x + e_\mu) = \omega(x) + a_\mu$, and therefore
\begin{gather*}
    (e^{i\alpha}U)^{u_t}(b) = e^{it\omega(x)}e^{i\alpha(b)}e^{i\zeta(b)}e^{-i(t\omega(x) + ta_\mu)} = e^{i\alpha_t(b)}e^{i\zeta(b)}\;,\\
    \alpha_t(b) \coloneqq \Phi\bigl(t\omega(x), \alpha(b), \zeta(b), ta_\mu\bigr)\;,
\end{gather*}
by the definition \eqref{eq:reduction-function}; the arguments have norm at most $\rho_0/2$ by hypothesis, and $\alpha_t$ has values in $\mathfrak g^c$ by (N2).

\emph{The plain differences of $\alpha$.} For $b, b + e$ in $B$, with $b = \langle x, x + e_\mu\rangle$, the value of $U$ on the bond from $x$ to $x + e$ is $e^{\pm i\zeta}$ with the value of $\zeta$ on that bond, so that by \eqref{eq:covariant-difference}, (N1), and \eqref{eq:reduction-hypotheses}
\[
    \abs{\alpha(b + e) - \alpha(b)} \leq \abs{(\delta_e\alpha)(b)} + c_Gz\,\abs{\alpha(b + e)} \leq a' + c_Gza = a''\;.
\]
Two bonds of the direction $\mu$ in $B$ are joined by a chain of at most $4D$ translations by unit vectors which stays in $B$, one coordinate after the other, because $B$ is a cube. Hence $\abs{\alpha(b) - a_\mu} \leq 4Da''$ for every bond $b$ of the direction $\mu$ in $B$.

\emph{The bounds \eqref{eq:reduced-phase}.} By (N2), $\alpha_t(b) = \alpha(b) - ta_\mu + \Psi(t\omega(x), \alpha(b), \zeta(b), ta_\mu)$, and by the first line of \eqref{eq:reduction-remainder}, with $\abs{t\omega} \leq 4aD$, the last term is at most $r$ in norm. For $t = 1$ the first two terms are at most $4Da''$ together, and for general $t$ at most $2a$.

\emph{The bound \eqref{eq:reduced-difference}.} By \eqref{eq:covariant-difference} and (N1), $\abs{(\delta_e\alpha_t)(b)} \leq \abs{\alpha_t(b + e) - \alpha_t(b)} + c_Gz(2a + r)$. The arguments of $\Phi$ at $b + e$ are $t\omega(x) + ta_e$, with $a_e \coloneqq \pm a_\nu$ for $e = \pm e_\nu$, then $\alpha(b + e)$, $\zeta(b + e)$, and the same $ta_\mu$. The difference of the linear parts is $\alpha(b + e) - \alpha(b)$, at most $a''$, and by the second line of \eqref{eq:reduction-remainder} the difference of the values of $\Psi$ is at most $c_G(2a + z)\abs{a_e} + c_G(4aD + 2a + z)a'' + c_G(4aD + 2a)z'$, the three terms being the changes of the first, the second, and the third argument, each multiplied by the bound of the corresponding derivative in \eqref{eq:reduction-remainder}.
\end{proof}

\subsection{One Carried Insertion}\label{sec:one-carried-insertion}

Let $W$ be a group of the deviation species made at the step $j$. We fix a point of the space of the step $k + 1$,
\begin{equation}\label{eq:carried-datum}
    (\mathbf U, \mathbf J) = (e^{i\alpha}U, \mathbf J) \in \mathcal S_{k+1}(U)\;,
\end{equation}
around a real configuration $U$ which is either $U_{k+1}$, for a block field $V_{k+1}$ in the support, with the current $J \coloneqq J_{k+1}$, or the trivial configuration, with $J \coloneqq 0$; in both cases $(U, J) \in \mathcal S_{k+1}(U)$ and, by (b), $U$ is one of the configurations of \eqref{eq:production-space}. By Lemma~\ref{lem:space-inclusion} and (a), both $(\mathbf U, \mathbf J)$ and $(U, J)$ lie in $\mathcal D_W$. We write
\begin{equation}\label{eq:carried-sizes}
\begin{gathered}
    a \coloneqq \alpha_{1,k+1}L^{-(k+1)}\;,
    \qquad
    \varrho_j \coloneqq \alpha_{1,j}L^{-j}\;,\\
    D_W \coloneqq 2u_\star M\sigma(W)L^j\;,
    \qquad
    \delta_W \coloneqq D_WL^{-(k+1)} = 2u_\star ML^{-m}\sigma(W)\;.
\end{gathered}
\end{equation}
By (c) the $\ell^\infty$-diameter of $X_W$ is at most $u_\star\sigma(W)ML^j$, and that of the dependence set $X^+_W$ exceeds it by at most $16M \leq ML^j$; hence $X^+_W$ lies in a cube $B_W$ of the fine lattice of side $D_W$, and $\delta_W$ is its side in the units of the scale $k + 1$. The interval has at most $D_W/2$ fine bonds, and $\delta_W \geq 2ML^{-m}$ because $\sigma(W) \geq 1$.

\begin{proposition}[One Carried Insertion]\label{prop:one-carried}
Let $\gamma$ be sufficiently small. With the notation above the following hold.
\begin{enumerate}[label=(\roman*)]
\item For every $W$,
\begin{equation}\label{eq:excess-plain}
    \norm{Y_W(\mathbf U, \mathbf J) - Y_W(U, J)} \leq c_{\mathrm{pl}}\,\varkappa_mL^{-m}\,\norm{Y_W}_{(j)}\;,
    \qquad
    c_{\mathrm{pl}} \coloneqq \max\Bigl(2,\ \frac{8c_G\lambda_\alpha}{c_FL}\Bigr)\;.
\end{equation}
\item Let $\mu \coloneqq \min\bigl(\frac12, c_FM/(c_G\lambda_\alpha)\bigr)$. If $\delta_W \leq \mu$, then
\begin{equation}\label{eq:excess-covariant}
    \norm{Y_W(\mathbf U, \mathbf J)} \leq 2\norm{Y_W(U, J)} + \varkappa_mL^{-m}\Bigl(\frac{40\,\delta_W}{\mu} + 14L^{-2m}\Bigr)e^{\nu\sigma(W)/8192}\,\norm{Y_W}_{(j)}\;.
\end{equation}
\end{enumerate}
\end{proposition}

The first bound is the gain of the radii alone and holds for every group; the second has the further factor $\delta_W \leq 2u_\star ML^{-m}\sigma(W)$ and is the one which is summable over the generations.

\begin{proof}
We write $r^{(n)} \coloneqq \alpha_{0,n}L^{-3n}$ for the radius of the current in \eqref{eq:space-of-step} and $a_2 \coloneqq \alpha_{1,k+1}L^{-2(k+1)} = aL^{-(k+1)}$ for that of the covariant differences at the step $k + 1$. By (d) and \eqref{eq:carried-blocking-condition},
\begin{equation}\label{eq:carried-ratios}
    \frac{a}{\varrho_j} \leq \varkappa_mL^{-m} \leq \frac{1}{20}\;,
    \qquad
    \frac{a_2}{\alpha_{1,j}L^{-2j}} \leq \varkappa_mL^{-2m}\;,
    \qquad
    \frac{r^{(k+1)}}{r^{(j)}} \leq \varkappa_mL^{-3m} \leq \frac{1}{20L^2}\;.
\end{equation}
The curvature of the real configuration is $f \coloneqq \max_p\abs{U(\partial p) - 1} \leq 2B_3\varepsilon_{k+1}L^{-2(k+1)}$ by (b), and $f = 0$ for the trivial configuration. Since $\varepsilon_{k+1}/\alpha_{0,k+1}$ and $\alpha_{1,n}$ tend to zero with $\gamma$, uniformly in the steps (Appendix~\ref{app:constants}, $q_0 > p_0$), we may and do assume that
\begin{equation}\label{eq:carried-smallness}
    c_Gf \leq \tfrac14c_F\,\alpha_{0,j}L^{-2j}\;,
    \qquad
    c_G\,\alpha_{1,n}^2 \leq \tfrac14c_F\,\alpha_{0,n}\;,
    \qquad
    \alpha_{1,n} \leq \rho_0
\end{equation}
for all steps $n$; the first uses $\alpha_{0,k+1} \leq \varkappa_m\alpha_{0,j}$ and $\varkappa_mL^{-2m} \leq 1$.

\emph{(i).} For $j \geq 2$ put
\[
    R' \coloneqq \min\Bigl(\frac{\varrho_j}{2a},\ \frac{c_F\,\alpha_{0,j}L^{-2j}}{8c_G\,a_2},\ \frac14\Bigl(\frac{r^{(j)}}{r^{(k+1)}} - 1\Bigr)\Bigr)\;,
\]
and for $j = 1$ let $R'$ be the first of the three numbers. If $2R' \leq 1$, the left side of \eqref{eq:excess-plain} is at most $2\norm{Y_W}_{(j)} \leq \norm{Y_W}_{(j)}/R'$, because both points lie in $\mathcal D_W$. Let $2R' > 1$. For $\abs z < 2R'$ the pair $(e^{iz\alpha}U, J + z(\mathbf J - J))$ lies in $\mathcal S_j(U)$: the complex part $z\alpha$ has norm less than $2R'a \leq \varrho_j$ and covariant differences less than $2R'a_2 \leq \varrho_jL^{-(k+1)} \leq \alpha_{1,j}L^{-2j}$; by \eqref{eq:complex-curvature} with $a_\circ = \varrho_j \leq \rho_0$ and $d_\circ = 2R'a_2$, and by \eqref{eq:carried-smallness} and the second entry of $R'$, its curvature is at most $c_Gf + 2c_GR'a_2 + c_G\varrho_j^2 \leq \frac34c_F\alpha_{0,j}L^{-2j}$; and the current has norm less than $r^{(k+1)} + 4R'r^{(k+1)} \leq r^{(j)}$, by the third entry. The function $g(z) \coloneqq Y_W(e^{iz\alpha}U, J + z(\mathbf J - J)) - Y_W(U, J)$ is therefore holomorphic on the disc of radius $2R' > 1$, it is at most $2\norm{Y_W}_{(j)}$ in norm there, and $g(0) = 0$, so that by Schwarz's lemma, applied to the matrix elements, $\norm{g(1)} \leq \norm{Y_W}_{(j)}/R'$. In both cases the left side of \eqref{eq:excess-plain} is at most $\norm{Y_W}_{(j)}/R'$, and by \eqref{eq:carried-ratios} the reciprocals of the three entries of $R'$ are at most $2\varkappa_mL^{-m}$, $(8c_G\lambda_\alpha/c_F)\varkappa_mL^{-2m}$, and $4 \cdot \frac{20}{19}\varkappa_mL^{-3m}$, each of which is at most $c_{\mathrm{pl}}\varkappa_mL^{-m}$.

\emph{(ii), the gauge.} The conditions \eqref{eq:space-of-step} are invariant under the $G$-valued gauge transformations, which conjugate $\alpha$ and $\mathbf J$, and so are the norms in \eqref{eq:excess-covariant}, by \eqref{eq:covariance-law} and because $\pi$ is unitary on $G$. For $U = U_{k+1}$ we apply Supplement, Lemma~17.22, to a cube $\mathcal Q$ of $2^4$ cubes of $\pi_{k+1}$ which contains $B_W$; such a cube exists because $\delta_W \leq \frac12$ is less than the side $M$ of these cubes. Its hypotheses hold for $\gamma$ small, because the local curvature is $E \leq 2B_3\varepsilon_{k+1}$ by (b) and $M_{\mathcal Q} \leq 4M$. It gives a $G$-valued gauge transformation after which $U = e^{i\zeta}$ on $B_W$ with, in fine units,
\[
    \abs{\zeta(b)} \leq z \coloneqq \hat bL^{-(k+1)}\;,
    \qquad
    \abs{\zeta(b + e) - \zeta(b)} \leq z' \coloneqq \hat bL^{-2(k+1)}\;,
    \qquad
    \hat b \coloneqq 24B_3^2M\varepsilon_{k+1}\;,
\]
and for the trivial configuration $\zeta = 0$. We work in this gauge and apply Lemma~\ref{lem:gauge-reduction} to the cube $B_W$, with $D = D_W$, $x_0 = x_W$, the radii $a$ and $a' = a_2$ of \eqref{eq:carried-datum}, and these $z$, $z'$; its hypotheses hold because $4aD_W = 4\alpha_{1,k+1}\delta_W \leq 2\alpha_{1,k+1}$ and $z$ are small. In the bounds of that lemma, $a_2D_W = a\delta_W$, $zaD_W = \hat b\,a\delta_W$, $a^2 = \alpha_{1,k+1}L^{-(k+1)}a \leq \alpha_{1,k+1}a\delta_W$, and $za = \hat b\,a_2$, so that every term is a multiple of $a\delta_W$, or of $a$, or of $a_2$, with a coefficient which is $4$, respectively $2$, respectively $1$, up to terms which vanish with $\alpha_{1,k+1}$ and $\hat b$. For $\gamma$ small we therefore have, for $t \in [0, 1]$ and on the bonds of $B_W$,
\begin{equation}\label{eq:reduced-datum}
    \abs{\alpha_1} \leq 5a\delta_W\;,
    \qquad
    \abs{\alpha_t} \leq 3a\;,
    \qquad
    \abs{\delta_e\alpha_t} \leq 2a_2\;,
    \qquad
    \abs\omega \leq 4\alpha_{1,k+1}\delta_W\;.
\end{equation}

\emph{(ii), the path.} We claim that $(\mathbf U^{u_t}, R(u_t)\mathbf J) \in \mathcal S_j(U)|_{X^+_W}$ for $t \in [0, 1]$. By Lemma~\ref{lem:gauge-reduction}, $\mathbf U^{u_t} = e^{i\alpha_t}U$, and by \eqref{eq:reduced-datum} and \eqref{eq:carried-ratios} the complex part satisfies $\abs{\alpha_t} \leq 3a < \varrho_j$ and $\abs{\delta_e\alpha_t} \leq 2a_2 \leq 2\varkappa_mL^{-2m}\alpha_{1,j}L^{-2j} < \alpha_{1,j}L^{-2j}$. For the curvature we use that $\delta_W \leq \mu$ and $\delta_W \geq 2ML^{-m}$ give $L^{-m} \leq \mu/(2M)$, so that $2a_2 \leq 2\varkappa_mL^{-m} \cdot L^{-m}\lambda_\alpha\alpha_{0,j}L^{-2j} \leq (\mu\lambda_\alpha/(20M))\alpha_{0,j}L^{-2j} \leq (c_F/(20c_G))\alpha_{0,j}L^{-2j}$ by the definition of $\mu$; with \eqref{eq:complex-curvature} for $a_\circ = 3a \leq \varrho_j$ and \eqref{eq:carried-smallness} the curvature is less than $c_F\alpha_{0,j}L^{-2j}$. Finally $\abs{R(u_t)\mathbf J} \leq (1 + c_G\abs\omega)\abs{\mathbf J} \leq 2r^{(k+1)} < r^{(j)}$ by (N1). For $j = 1$ only the first of these conditions is posed. Since $u_1(x_W) = 1$, \eqref{eq:base-point-invariance} and \eqref{eq:base-point-form} give
\begin{equation}\label{eq:carried-identity}
    Y_W(\mathbf U, \mathbf J) = \hat Y_W(\mathbf U^{u_1}, R(u_1)\mathbf J)\,\pi\bigl(\mathbf U(I_W)\bigr)\;.
\end{equation}

\emph{(ii), the disc.} Put
\[
    R_A \coloneqq \frac{\mu\varrho_j}{10\,a\delta_W}\;,
    \qquad
    R_J \coloneqq \frac13\Bigl(\frac{r^{(j)}}{r^{(k+1)}} - 1\Bigr)\;,
    \qquad
    R \coloneqq \min(R_A, R_J)\;,
\]
with $R \coloneqq R_A$ for $j = 1$, and $f(z) \coloneqq \hat Y_W(e^{iz\alpha_1}U, J + z(R(u_1)\mathbf J - J))$ for $\abs z < R$, so that $f(1)$ is the first factor of \eqref{eq:carried-identity} and $f(0) = \hat Y_W(U, J)$. Since $\delta_W \leq \mu$, we have $R_A \geq \varrho_j/(10a) \geq 2$ and $R_J \geq 2$ by \eqref{eq:carried-ratios}. For $\abs z < R$ the pair in the argument of $f$ lies in $\mathcal S_j(U)|_{X^+_W}$: by \eqref{eq:reduced-datum} the complex part $z\alpha_1$ has norm less than $5R_Aa\delta_W = \mu\varrho_j/2$; its covariant differences are less than $2R_Aa_2 = \mu\varrho_jL^{-(k+1)}/(5\delta_W) \leq \mu\,\alpha_{1,j}L^{-2j}/(10M)$, because $\delta_W \geq 2ML^{-m}$; by \eqref{eq:complex-curvature}, \eqref{eq:carried-smallness}, and the definition of $\mu$ its curvature is at most
\[
    c_Gf + \frac{c_G\mu\lambda_\alpha}{10M}\alpha_{0,j}L^{-2j} + c_G\varrho_j^2 \leq \Bigl(\frac14 + \frac{1}{10} + \frac14\Bigr)c_F\,\alpha_{0,j}L^{-2j}\;;
\]
and the current has norm less than $r^{(k+1)} + 3R_Jr^{(k+1)} = r^{(j)}$, because $\abs{R(u_1)\mathbf J} \leq 2r^{(k+1)}$. Hence $f$ is holomorphic on the disc of radius $R$. Along the interval the transport of $e^{iz\alpha_1}U$ is a product of at most $D_W/2$ factors $e^{iz\alpha_1(b)}U(b)$, with $\pi(U(b))$ unitary and $\norm{\pi(e^{iX})} \leq e^{c'_\pi\abs X}$ (Section~\ref{sec:localisation-expansion}), so that by \eqref{eq:base-point-form} and (a)
\begin{align*}
    \norm{f(z)} &\leq \norm{Y_W}_{(j)}\exp\bigl(c'_\pi\tfrac12\mu\varrho_j \cdot \tfrac12D_W\bigr) = \norm{Y_W}_{(j)}\exp\bigl(\tfrac12c'_\pi\mu u_\star M\alpha_{1,j}\,\sigma(W)\bigr)\\
    &\leq e^{\nu\sigma(W)/8192}\,\norm{Y_W}_{(j)}\;,
\end{align*}
the last for $\gamma$ so small that $4096\,c'_\pi u_\star M\alpha_{1,n} \leq \nu$ for all $n$. By Schwarz's lemma for $f - f(0)$, which is at most twice this bound on the disc, $\norm{f(1) - f(0)} \leq (2/R)e^{\nu\sigma(W)/8192}\norm{Y_W}_{(j)}$. In the same way $\norm{\pi(\mathbf U(I_W))} \leq \exp(c'_\pi a D_W/2) = \exp(\frac12c'_\pi\alpha_{1,k+1}\delta_W) \leq 2$, and $\norm{f(0)} = \norm{Y_W(U, J)}$ because $\pi(U(I_W))$ is unitary. By \eqref{eq:carried-identity},
\[
    \norm{Y_W(\mathbf U, \mathbf J)} \leq 2\bigl(\norm{f(0)} + \norm{f(1) - f(0)}\bigr) \leq 2\norm{Y_W(U, J)} + \frac4R\,e^{\nu\sigma(W)/8192}\norm{Y_W}_{(j)}\;,
\]
and $4/R \leq 4/R_A + 4/R_J \leq 40\varkappa_mL^{-m}\delta_W/\mu + 12 \cdot \frac{20}{19}\varkappa_mL^{-3m}$ by \eqref{eq:carried-ratios}, which is \eqref{eq:excess-covariant}.
\end{proof}

\subsection{The Sum over the Generations}\label{sec:generation-sum}

\begin{proposition}[The Carried Groups in the Supremum Clause]\label{prop:carried-sup}
Let the ensemble be dressed at the steps $1, \dots, k$, and suppose that on the term without large-field regions of the step $k + 1$ the groups carried from these steps satisfy \eqref{eq:per-insertion} and \eqref{eq:tail-clause} at the step $k + 1$, at every real background of the term, and vanish at its reference data. Let $\gamma$ be sufficiently small and $L > \Lambda_{\mathfrak a} \coloneqq (1 + \beta_0)^{10}$. Then, for every traversal $\mathfrak n'$ of the step $k + 1$,
\begin{equation}\label{eq:carried-sup}
    \sum_{W \ni \mathfrak n'}e^{\theta(W) + \lambda_c\nu\sigma(W) + \nu_e\operatorname{ext}(W)}\norm{Y_W}_{(k+1)} \leq 2C_\chi\tilde\rho_{k+1}\bar\delta^*_{\mathfrak n'} + \hat C\,C_{S'}\,\mathfrak a_{k+1}\;,
\end{equation}
the sum over the groups of the deviation species made at the steps $j \leq k$ with $\mathfrak n'$ in the interval, where $\bar\delta^*_{\mathfrak n'}$ is the expression \eqref{eq:deviation-quantity} of the step $k + 1$ with each maximum replaced by its supremum over the real backgrounds of the term, and
\begin{equation}\label{eq:generation-constant}
    \hat C \coloneqq \sum_{m \geq 1}\Lambda_{\mathfrak a}^m\,\varkappa_m\Bigl[L^{-m}\Bigl(\frac{1280\,u_\star M}{e\,\nu\mu} + 14L^{-m}\Bigr) + c_{\mathrm{pl}}\exp\Bigl(-\frac{\nu\mu L^m}{32u_\star M}\Bigr)\Bigr]\;.
\end{equation}
Here $\theta(W)$, $\sigma(W)$, and $\operatorname{ext}(W)$ are the numbers fixed when $W$ was made, as in \eqref{eq:production-clause}.
\end{proposition}

The constant $\hat C$ is finite because $L > \Lambda_{\mathfrak a}$, and it depends on $L$, $M$, $\nu$, $\beta_0$, $u_\star$, $c_F$, $c_G$, and $\lambda_\alpha$, and not on $k$ or on the coupling. The first term of \eqref{eq:carried-sup} is at most $2C_\chi\tilde\rho_{k+1} \cdot 2c_\sigma B_3\varepsilon_{k+1}W_1$, because the bound \eqref{eq:deviation-scaling} holds for each maximum and at every real background, and this is less than $\mathfrak a_{k+1}$ for $\gamma$ small, since $\tilde\rho_{k+1}\varepsilon_{k+1} = \rho_{k+1}\varepsilon_{k+1}/\min(\alpha_{0,k+1}, \alpha_{1,k+1})$ and $q_0 > p_0$. With \eqref{eq:production-clause} at the step $k + 1$, in which $\ell_W \geq 1$, $S_W \geq \norm{Y_W}_{(k+1)}$, and $\lambda_c \leq 1$, the proposition therefore gives \eqref{eq:sup-clause} at the step $k + 1$ with
\begin{equation}\label{eq:sup-constant}
    C_c \coloneqq 1 + (1 + \hat C)\,C_{S'}\;.
\end{equation}
The constant does not come from a recursion: the bound of the carried groups uses the production clauses of the earlier steps and no supremum clause.

\begin{proof}
\emph{The amplitudes of two steps.} We first show that $\mathfrak a_n \leq \Lambda_{\mathfrak a}\mathfrak a_{n+1}$ for $1 \leq n \leq k$, so that $\mathfrak a_j \leq \Lambda_{\mathfrak a}^m\mathfrak a_{k+1}$. By \eqref{eq:tube-two-six} and \eqref{eq:tube-two-seven}, $g_n \leq (1 + \beta_0)g_{n+1}$ and, for $\gamma$ small, $(1 + \beta_0)^{-1}r_{n+1} \leq r_n \leq (1 + \beta_0)r_{n+1}$ for the radii $r_n = A_1(\log g_n^{-2})^{p_0}$ of \eqref{eq:class-parameters}. Hence $\rho_n \leq (1 + \beta_0)^6\rho_{n+1}$ and $g_n(1 + r_n)^3 \leq (1 + \beta_0)^4g_{n+1}(1 + r_{n+1})^3$, while $e^{-c_\tau r_n^2} \leq e^{-c_\tau r_{n+1}^2/(1 + \beta_0)^2} \leq \beta_0g_{n+1}$ for $\gamma$ small, because $r_{n+1}^2$ is a multiple of the power $2p_0 > 1$ of $\log g_{n+1}^{-2}$. Therefore $\tau_n \leq ((1 + \beta_0)^4 + \beta_0)\tau_{n+1} \leq (1 + \beta_0)^5\tau_{n+1}$, and the claim follows from \eqref{eq:amplitude}.

\emph{The real points.} Let $W$ be a group of the sum, made at the step $j$, and \eqref{eq:carried-datum} a point of the space of the step $k + 1$, around $U$. We apply Proposition~\ref{prop:one-carried}(ii) if $\delta_W \leq \mu$ and (i) otherwise, and in both cases
\[
    \norm{Y_W(\mathbf U, \mathbf J)} \leq 2\norm{Y_W(U, J)} + \epsilon_W\;,
\]
with $\epsilon_W$ the second term of \eqref{eq:excess-covariant}, respectively the right side of \eqref{eq:excess-plain}; it does not depend on the point. If $U$ is the trivial configuration, the first term vanishes by hypothesis. If $U = U_{k+1}$, let $\mathfrak n_j \subset \mathfrak n'$ be a traversal of the step $j$ in the interval of $W$ and $t \coloneqq t^{(k+1)}_{\mathfrak n'}(W)$. As in the proof of Lemma~\ref{lem:tail-to-sum}, the comparison region of $W$ at $\mathfrak n_j$ lies within the distance $R_t$ of $\mathfrak n'$ in the units of the scale $k + 1$, so that by \eqref{eq:per-insertion} and \eqref{eq:insertion-deviation} $e^{\theta(W)}\norm{Y_W(U, J)} \leq \beta_W\bar m_t$, where $\bar m_t$ is the supremum over the real backgrounds of the maximum of $\abs{U(\partial p) - 1}$ over $p \in \mathcal P_t(\mathfrak n')$. Taking the supremum over the points we obtain
\begin{equation}\label{eq:carried-pointwise}
    e^{\theta(W)}\norm{Y_W}_{(k+1)} \leq 2\beta_W\,\bar m_{t^{(k+1)}_{\mathfrak n'}(W)} + e^{\theta(W)}\epsilon_W\;.
\end{equation}
The numbers $\bar m_t$ are nondecreasing in $t$, and the summation by parts of the proof of Lemma~\ref{lem:tail-to-sum}, with \eqref{eq:tail-clause} at the step $k + 1$ and $\lambda_c \leq 1$, bounds the weighted sum of the first terms of \eqref{eq:carried-pointwise} over the groups of the sum by $2C_\chi\tilde\rho_{k+1}\bar\delta^*_{\mathfrak n'}$.

\emph{The excess of one generation.} Fix $j \leq k$ and consider the groups of the sum which were made at the step $j$. Each of them has a traversal of the step $j$ in its interval which is contained in $\mathfrak n'$, namely the one which contains a fine bond common to $\mathfrak n'$ and the interval, and $\mathfrak n'$ contains $L^m$ traversals of the step $j$. By \eqref{eq:production-clause} at the step $j$, with $\ell_W \geq 1$ and $S_W \geq \norm{Y_W}_{(j)}$,
\begin{equation}\label{eq:generation-production}
    \sum_We^{\theta(W) + \nu\sigma(W) + \nu_e\operatorname{ext}(W)}\norm{Y_W}_{(j)} \leq L^mC_{S'}\,\mathfrak a_j\;,
\end{equation}
the sum over these groups. For those with $\delta_W \leq \mu$ we have $\lambda_c\nu + \nu/8192 = \frac{15}{16}\nu$ and, by \eqref{eq:carried-sizes} and $te^{-\nu t/16} \leq 16/(e\nu)$,
\begin{align*}
    e^{\lambda_c\nu\sigma(W)}\epsilon_W &\leq e^{\nu\sigma(W)}\varkappa_mL^{-m}\Bigl(\frac{80u_\star ML^{-m}}{\mu}\sigma(W)e^{-\nu\sigma(W)/16} + 14L^{-2m}\Bigr)\norm{Y_W}_{(j)}\\
    &\leq e^{\nu\sigma(W)}\varkappa_mL^{-2m}\Bigl(\frac{1280\,u_\star M}{e\,\nu\mu} + 14L^{-m}\Bigr)\norm{Y_W}_{(j)}\;.
\end{align*}
For those with $\delta_W > \mu$ we have $\sigma(W) > \mu L^m/(2u_\star M)$, so that $e^{\lambda_c\nu\sigma(W)} \leq e^{\nu\sigma(W)}e^{-\nu\sigma(W)/16} \leq e^{\nu\sigma(W)}\exp(-\nu\mu L^m/(32u_\star M))$ and
\[
    e^{\lambda_c\nu\sigma(W)}\epsilon_W \leq e^{\nu\sigma(W)}c_{\mathrm{pl}}\,\varkappa_mL^{-m}\exp\Bigl(-\frac{\nu\mu L^m}{32u_\star M}\Bigr)\norm{Y_W}_{(j)}\;.
\]
Multiplying by $e^{\theta(W) + \nu_e\operatorname{ext}(W)}$ and summing with \eqref{eq:generation-production}, we find that the weighted sum of the $\epsilon_W$ over the groups made at the step $j$ is at most the $m$-th term of \eqref{eq:generation-constant} without the factor $\Lambda_{\mathfrak a}^m$, times $C_{S'}\mathfrak a_j$, and $\mathfrak a_j \leq \Lambda_{\mathfrak a}^m\mathfrak a_{k+1}$.

\emph{The sum.} Summing the excesses over $j \leq k$, that is, over $m \geq 1$, and adding the first terms of \eqref{eq:carried-pointwise}, we obtain \eqref{eq:carried-sup}.
\end{proof}

\bibliographystyle{amsplain}
\bibliography{refs}

\end{document}
